\section{Proposed framework}
\label{sec:method}

The framework has three components. The referee randomises its accept-or-reject decisions and
logs their propensities, which turns the controller's logs into an experiment with a known design. The predictable
matched-filter randomisation test (PMRT) uses this design to test each hypothesis of Section~\ref{sec:estimand} and
returns a signed causal effect map with asymptotic control of the FDR. A map-driven referee then uses the map to
reject harmful requests, and an edge-certification gate restricts it to rejections whose side effects on the guard
KPIs are certified small.

\subsection{Test statistic}
\label{sec:pmrt}\claim{DEF-PMRT-WEIGHT}

In reactive logs, network state drives both the mode $W$ and the KPIs. If the referee accepts when an independent
uniform draw $U$ falls below a logged propensity $\pi$, state affects the mode only through $\pi$. PMRT tests whether a
KPI moves with the modes more than those propensities predict, pairing a score of each mode with a weight built from
the KPI series and comparing the statistic with redraws from the logged design. It establishes whether an effect
exists, not its size.

Take a family $f$ and its units $u$ in information order: episode, opening time, cell. With $L(\text{accept}) = 1$ and
$L(\text{reject}) = 0$, the assignment score of unit $u$ is
\begin{equation}
v_u = \mathrm{sgn}_u \bigl( L(W_u) - \pi_u \bigr),
\label{eq:score}\claim{DEF-PMRT-SCORE}
\end{equation}
where $\mathrm{sgn}_u$ is the direction of the opening request. Over the policy's draw, $v_u$ has mean zero and
variance $\pi_u(1 - \pi_u)$, so an acceptance that was likely scores little and one against the odds scores more.
The weight $w_u$ summarises how the KPI under test moved after the unit opened. It is the residual of the KPI after a
ridge prediction~\cite{hoerl1970ridge} from what was known at opening, combined over the receiving cells and the
seconds after opening by a matched filter, then centred, shrunk and clipped~\cite{huber1964robust}. Every fitted part
was trained on earlier episodes, separate from the evaluation data (\esm{esm:method}).

The validity of the test rests on one property of the weight: apart from the KPI series under test, it uses only
information fixed before the unit drew its mode. It therefore uses no later unit, no unit of another family opening at
or after $u$, and no mean or slope over an episode's realised units, because whether those units exist can depend on
$W_u$. The statistic and its variance are
\begin{equation}
S = \sum_u v_u w_u, \qquad
\sigma^2 = \sum_u \pi_u (1 - \pi_u)\, w_u^2, \qquad
z = S / \sigma,
\label{eq:stat}\claim{DEF-PMRT-STAT}
\end{equation}
where $\sigma^2$ is the variance of $S$ over redraws of the modes with the weights held. Since it is the same for every
redraw, $z$ is normalised by the design variance rather than studentised. The reported sign is that of $S$ times an
orientation $o$, the sign of the training effect, which is set only when that effect is clear in the training data.

\subsection{Inference and false discovery rate control}
\label{sec:inference}\claim{DEF-PMRT-SUPPORT}\claim{DEF-FDR-WBY}\claim{ASM-A1-DESIGN}\claim{ASM-A2-NULL}\claim{ASM-A3-PRED}\claim{ASM-A4-REG}\claim{CIT-BY}\claim{CIT-RW}

The modes of family $f$ are redrawn $B$ times from the logged propensities, with every other family fixed and
the weights held, and $z_b$ is recomputed for each redraw. Each hypothesis has a two-sided p-value $p$ and, where $o$
is set, a one-sided value $p_d$ for each effect direction $d \in \{+1, -1\}$:
\begin{equation}
p = \frac{1 + \#\{b : |z_b| \ge |z|\}}{B + 1}, \qquad
p_d = \frac{1 + \#\{b : d\, o\, z_b \ge d\, o\, z\}}{B + 1},
\label{eq:pval}\claim{DEF-PMRT-PVAL}
\end{equation}
where $\#$ counts redraws. Algorithm~\ref{alg:pmrt} summarises the declaration pipeline; the complete per-hypothesis
algorithm and fixed support, window and weight-construction parameters are given in \esm{esm:method}. We use
$B = 9999$ redraws.

\begin{algorithm}[tbp]
\caption{PMRT declarations.}
\label{alg:pmrt}
\footnotesize
\begin{algorithmic}[1]
\Require hypotheses $\mathcal{H}$; training and evaluation units with modes, request directions, propensities and KPI
  series; redraws $B$
\Ensure signed declarations or untested hypotheses
\For{each hypothesis $j = (f, r, k) \in \mathcal{H}$}
  \State Check the required units, accepted and rejected modes, episodes and outcome windows; otherwise mark $j$
    untested
  \State On training episodes, fit the pre-window adjustment, matched filter, scale, clipping rule and prior direction
  \State For each evaluation unit in information order, compute its assignment score $v_u$ and outcome weight $w_u$
  \State Compute the observed statistic $z$ from Eq.~\eqref{eq:stat}
  \State Redraw the modes of family $f$ from their logged propensities $B$ times, holding the weights and other
    families fixed
  \State Compute the two-sided and available directional p-values from Eq.~\eqref{eq:pval}, with the reported sign
\EndFor
\State Apply the weighted BY step-up to all hypotheses; untested hypotheses are never declared
\State \Return the rejected hypotheses and their reported signs
\end{algorithmic}
\end{algorithm}

The sharp null $H_0(f, r, k)$ of Section~\ref{sec:estimand} is stated over possible units, because which units are
realised depends on the modes. The history $\mathcal{G}_{u-1}$ holds all other randomness and the modes of $f$'s first
$u - 1$ units, and the result rests on four assumptions.

\begin{description}
\item[(A1) Design.] For every unit $u$, its existence, cell, opening time, direction $\mathrm{sgn}_u$ and propensity
  $\pi_u$ are $\mathcal{G}_{u-1}$-measurable, $\Pr(W_u = \text{accept} \mid \mathcal{G}_{u-1}) = \pi_u$, and
  $\pi_{\min} \le \pi_u \le 1 - \pi_{\min}$ for a fixed $\pi_{\min} > 0$.
\item[(A2) Sharp null.] $H_0(f, r, k)$ holds.
\item[(A3) Predictable weights.] Every $w_u$ is $\mathcal{G}_{u-1}$-measurable and square integrable.
\item[(A4) Regularity.] As the number of episodes grows, the design variance $\sigma^2$ grows without any single unit
  dominating $S$.
\end{description}

Operationally, the RIC must draw each mode from its logged, past-determined propensity bounded away from 0 and 1; the
analyst must build each weight from information available when the unit opened, apart from the KPI under test; and the
log must contain enough units that none dominates $S$. In the simulator the last condition holds when clipping keeps
the weights bounded and $\sigma^2/n$ approaches a positive limit.

\begin{proposition}[Asymptotic validity of PMRT]
\label{prop:pmrt}\claim{ASM-ORIENT}\claim{THM-PMRT-ASY}
Assume (A1) to (A3). Then $S = \sum_u v_u w_u$ is a martingale with respect to $(\mathcal{G}_u)$, and its predictable
quadratic variation is $\sigma^2$. Let (A4) also hold and the number of episodes grow.
\begin{enumerate}
\item The statistic $z = S / \sigma$ converges in distribution to $\mathcal{N}(0, 1)$.
\item Conditionally on the data, every redraw $z_b$ has mean 0 and, where $\sigma > 0$, variance 1, and its distribution
  function converges uniformly to $\Phi$ in probability.
\item For every fixed $B$ and every $t \in [0, 1]$, $\limsup \Pr(p \le t) \le \lfloor t(B + 1) \rfloor / (B + 1) \le t$.
  The same holds for $p_d$ for each $d$ when $o$ is fixed in advance.
\end{enumerate}
The p-values $p$ and $p_d$ are therefore asymptotically valid. The value $p$ differs from
$2\{1 - \Phi(|z|)\}$ only by a term that vanishes in probability and by Monte Carlo error.
\end{proposition}

Adaptive unit creation, repeated requests and state-dependent propensities do not invalidate the martingale argument
when their existence, timing and propensity are determined before the current mode is drawn, as (A1) requires. The
fresh-draw lemma\claim{THM-FRESH} (\esm{lem:s_fresh}) then gives each assignment score conditional mean zero, so under
the sharp null $S$ has predictable variance $\sigma^2$ in every sample. Assumption (A4) supplies the asymptotic
calibration through the martingale central limit theorem~\cite[Cor.~3.1 and 3.2]{hall1980martingale}\claim{CIT-HH} and
the Lindeberg--Feller and P\'olya theorems~\cite[Thm.~3.4.10 and Exercise~3.2.9]{durrett2019probability}\claim{CIT-DUR}.
The redraw p-value is nevertheless not exact in finite samples: changing a mode can alter the later network state,
units and propensities, so the adaptive design cannot generally be replayed as required by exact sequential
randomisation tests~\cite{bojinov2019timeseries}. A fixed design gives a valid tie-inclusive p-value\claim{THM-EXACT},
exact for a continuous statistic~\cite[Sect.~4]{phipson2010permutation}\claim{CIT-PS}, but PMRT does not meet that
condition; a two-unit counterexample is given in \esm{ex:s_adaptive}\claim{THM-COUNTEREX}. The placebo and null-outcome
experiments of Section~\ref{sec:res_validity} assess finite-sample behaviour.

The 60 hypotheses are declared together by a weighted Benjamini--Yekutieli (BY)
step-up~\cite[Thm.~1.3]{benjamini2001dependency}, \cite{blanchard2008fdr} at $q = 0.05$. Hypothesis $j$ enters with
training-derived weight $\omega_j$ and, when available, prior-direction p-value $p_{d_j}$; otherwise it enters with
$p$~\cite[Thm.~3.1 and Sect.~5]{roeder2009weighted}. A rejection is reported with that prior direction or the test's
reported sign.

\begin{corollary}[Asymptotic false discovery rate]
\label{cor:fdr}\claim{ASM-FDR-FIXED}\claim{THM-FDR-ASY}
Let $\mathcal{H}_0$ hold the true sharp nulls among the $m = 60$ hypotheses. Let (A1) to (A4) hold for each of them,
and let the weighted step-up use directions $d_j$ and weights $\omega_{n,j} \ge 0$ with $\sum_j \omega_{n,j} = m$ for
the analysed sample size $n$, all fixed from data independent of the evaluation data, with $\omega_{n,j} \to \omega_j$
(or $\omega_{n,j} = \omega_j$ for every $n$). Then, whatever the dependence between the p-values,
\begin{equation}
\limsup_{n \to \infty} \mathrm{FDR} \le \frac{q}{m} \sum_{j \in \mathcal{H}_0} \omega_j \le q,
\label{eq:fdr}
\end{equation}
where a false discovery is the rejection of a true sharp null.
\end{corollary}

The result applies the reshaped weighted step-up of Blanchard and Roquain~\cite[Thm.~4.2]{blanchard2008fdr}
\claim{CIT-BR} through a finite-sample bound with asymptotically vanishing slack\claim{THM-BY-FIN}
(\esm{esm:theory}). Counting untested hypotheses in $m$ is conservative because they are never declared.

\subsection{Scope of the guarantees}
\label{sec:reading}\claim{THM-SCOPE}

A declared edge states that the modes of family $f$ change KPI $k$ somewhere in the cells that relation $r$ covers,
and the FDR bound covers that claim asymptotically. It does not cover the reported sign, effect magnitude or unique
spatial attribution: a unit's neighbour and far cells are the own cells of other units of $f$, so an effect of $f$ on
its own cells also moves the KPI of their neighbour and far cells, and the far hypotheses are no negative control
(\esm{tab:s_attr}). Table~\ref{tab:terms} summarises these distinctions. The gate similarly screens undeclared guard
edges but does not certify the resulting policy, whose constraints are evaluated on held-out seeds.

\subsection{Map-driven referee}
\label{sec:slope}\claim{ASM-SLOPE-CENTRE}
\label{sec:referee}\claim{DEF-REF-RULES}

The referee needs the size of each effect, not only its sign, and takes it from the design slope $\hat\beta$, which
regresses each unit's centred outcome on its assignment score. Let $e_u = y_u - \tilde y_u$ be the outcome minus an
outcome centre $\tilde y_u$ computed from earlier units' outcomes and $u$'s pre-window. Under the uniform-randomised policy,
\begin{equation}
\hat\beta = \frac{\sum_u v_u e_u}{\sum_u v_u^2}
          = \frac{1}{N} \sum_u \mathrm{sgn}_u \Bigl( \frac{\mathbf{1}\{W_u = \text{accept}\}\, e_u}{\pi_u}
            - \frac{\mathbf{1}\{W_u = \text{reject}\}\, e_u}{1 - \pi_u} \Bigr),
\label{eq:slope_ht}\claim{DEF-SLOPE}
\end{equation}
with $\pi_u = 0.5$, which is the Horvitz--Thompson estimator of $\tau_{\pi}$ and converges to it
(design-slope lemma, \claim{THM-SLOPE}\esm{esm:theory}). Under a state-dependent policy such as the incumbent, it converges to the
$\pi_u(1 - \pi_u)$-weighted average of $\mathrm{sgn}_u \tau_u^{\pi}$. The lemma assumes an outcome centre fixed before
$u$'s mode is drawn, which the implemented centre meets only approximately (Section~\ref{sec:disc_limits}).

A causal effect map $M$ gives each declared edge $(f, r, k)$ a signed coefficient $\beta(f, r, k)$. For PMRT it is the
declaration's sign times the magnitude of the design slope; an associational map uses the method's sign times the
magnitude of the least-squares slope of the outcome on the logged mode. The referee decides when a decision unit opens,
on the opening request's family, direction and context, and later requests of the pair inherit the mode. For a request
of family $f$ and direction $d$ it sums the direction-adjusted effects $\Delta_k = d \sum_r \beta(f, r, k)$ over the
own, neighbour and far relations for each KPI $k$, judges the request on the first KPI with $\Delta_k \ne 0$ in the
order protected-slice violations, service violations, RLFs, and rejects it if it is harmful and either costly in energy
or under protected-UE pressure. A rejection that accepting would have helped on a guard KPI is delayed by at most one
unit. The same constants apply to every map; the full rules are in \esm{esm:method}.

\subsection{Edge-certification gate}
\label{sec:gate}\claim{DEF-GATE}\claim{ASM-GATE-COV}\claim{THM-GATE-BUDGET}

A map leaves some guard edges undeclared, and a referee acting on it can harm a guard KPI through such an edge. The
gate prevents this by letting the referee reject a class of requests only when every undeclared guard edge of the
class is certified small. It checks every class $(f, d)$ that the referee can reject against the own, neighbour and far
relations of the two guard KPIs, service violations and RLFs. Rejecting a request of direction $d$ moves a guard KPI
by $-d\,\beta$, so a rejection's harmful direction is $h = -d$. We write the class $(f, +1)$ of requests that raise $f$
as $f$-up, so that sleep-up holds the requests to put a pico to sleep, and $(f, -1)$ as $f$-down.

An undeclared edge is certified when the one-sided 90~\% upper bound of $h\,\beta$ from the map's own estimator is at
most a tolerance $\varepsilon(f, d, k)$, and is otherwise unresolved; a class with any unresolved edge is always
accepted, and a map without an effect estimator cannot certify. The tolerance is
\begin{equation}
\varepsilon(f, d, k) = \frac{0.05\, G_k}{3\, \max\{N(f, d), 1\}},
\label{eq:tol}
\end{equation}
with $G_k$ the mean per-episode total of guard KPI $k$ and $N(f, d)$ the mean number of units of the class per
episode, both measured under accept-all on development episodes. Suppose that every unit of a class were rejected,
that the class had at most $\max\{N(f, d), 1\}$ units, and that each rejection moved the KPI by at most $\varepsilon$ on
each of the three relations. The guard ratio would then move by at most 0.05, half the margin of 0.10 that a referee
may use under the 1.10 limit of Section~\ref{sec:estimand}. The certificate is thus a first-order budget per edge, not
a bound on the referee's guard ratios: it ignores overlapping windows, classes acting together and declared edges, so
whether a referee meets its guards is checked on held-out seeds.
